% !TeX program = xelatex
% 该文件仅演示模板如何填写，不代表最终课程内容。

\documentclass{hepingjie_lessonplan}

\lessonplansetup{
  topic         = {数轴},
  type          = {新授课},
  chapter-total = {8},
  topic-hours   = {1},
  lesson-no     = {3},
  year          = {2026},
  month         = {9},
  day           = {10},
  record-page   = {1}
}

\begin{document}

\MakeLessonOverview
  {理解数轴的概念，掌握原点、正方向和单位长度三个要素；能够在数轴上表示有理数，并依据点的位置比较简单有理数的大小。}
  {数轴三要素；在数轴上准确表示有理数。}
  {单位长度的统一以及负数在数轴上的定位。}
  {情境引入、观察归纳、数形结合、练习反馈。}
  {多媒体课件、直尺、学生学案。}
  {左侧板书数轴三要素；中部保留标准数轴；右侧板书例题与易错提醒。}

\MakeTeachingProcessPage
  {\LessonStage{一、情境导入}{借助温度计和道路里程标，引导学生思考：怎样用一条直线表示正数、负数和零？}
   \LessonStage{二、形成概念}{教师示范画一条直线，依次确定原点、正方向和单位长度，给出数轴定义。}
   \LessonStage{三、辨析三要素}{展示若干“数轴”，让学生判断是否规范，并说明理由。}}
  {观察情境；回答问题；在学案上补画原点、箭头和刻度。}
  {5分钟\\12分钟\\8分钟}

\MakeTeachingProcessPage
  {\LessonStage{四、表示有理数}{示范在数轴上表示 $3$、$-2$、$\frac12$、$-3.5$，强调先判断方向，再确定距离。}
   \LessonStage{五、课堂练习}{学生独立完成基础题，教师巡视并选取典型错误讲评。}
   \LessonStage{六、归纳提升}{归纳“数轴上的点与有理数”的对应关系，并渗透数形结合思想。}}
  {独立作图；同桌互查；订正错误；口头归纳。}
  {10分钟\\7分钟\\3分钟}

\MakeTeachingProcessReflectionPage
  {\LessonStage{七、课堂小结}{学生从“知识、方法、易错点”三个角度回顾本节课。}
   \LessonStage{八、作业布置}{完成教材本节练习，并订正课堂学案。}}
  {完成自评；记录作业。}
  {3分钟\\2分钟}
  {记录课堂中学生对“单位长度统一”的掌握情况；关注数轴刻度不均、负数方向错误等问题，并据此调整下一课时的复习安排。}

\end{document}
